% Korean edition of figs/pytimings.tex.
\fig[wide]{pytimings}{Python 클라이언트는 프로세스당 한 번 1.2\textasciitilde 1.3\,s, 재사용 요청마다 31\textasciitilde 39\ms{}를 쓴다}[클라이언트, 시점, 연결 상태별 REST 요청 하나의 시간, 2026년 11월 24일]{%
{\tabfont\begin{tabularx}{\linewidth}{>{\RaggedRight\arraybackslash}p{40mm} >{\RaggedRight\arraybackslash}p{42mm} R{17mm} @{\hspace{6pt}}>{\RaggedRight\arraybackslash}X}
\thd{클라이언트, 시점} & \thd{요청} & \thd{시간} & \thd{시간이 드는 곳}\\
curl, 06:03 (개장 전) & 재사용 연결 & 16.0\ms & 왕복 1회(11.1\ms)와 서버에서 약 5\ms{} (\figref{timing})\\\rowsep
\hl Python, 장중 & 재사용 연결 & 31\textasciitilde 39\ms & curl과의 차이 원인은 분리하지 않음: 클라이언트, 시각, 장중 부하\\\rowsep
Python, 새 프로세스 & 첫 요청 & 1.2\textasciitilde 1.3\,s & 모듈 로딩과 SSL 초기화, 프로세스당 한 번. 가상환경이 Windows 드라이브에 있음\\\rowsep
Python, 같은 프로세스 & 두 번째 클라이언트의 첫 요청 & 130\ms & 새 연결 비용만. 일회성 비용은 이미 치름\\
\end{tabularx}}}{음영 행은 시작 이후 래퍼 요청마다 드는 비용이다. 1.2\textasciitilde 1.3\,s는 시작할 때 첫 주문 전에 \code{warm_up()}이 한 번 치른다. Python 시간은 \code{transport} 실서버 검사(\figref{tests})에서 나왔다.}{DART Quant Team}
